\section{未锚定账本事件的叙事补充：western-xia}

\label{sec:anchor-batch-two-western-xia}

本组为尚无正文跳转目标的账本事件补入锚点。每一记录置于战争、行政、环境、迁徙与地方社会的关系中，并按来源限度约束叙事。

\subsection{制度、区域与材料边界}

\eventanchor{XIA-1053-REGNAL-YEAR}{1053年·西夏以李谅祚在位第5年作为诏令、赋役与外交文书的纪年框架}账本条目记录西夏在1053年的“西夏以李谅祚在位第5年作为诏令、赋役与外交文书的纪年框架”。其背景是新岁颁历、年号和纪年是中枢与地方行政沟通的共同前提，后续影响为为同年政令、战事和地方材料提供日期锚点，不能单独推知具体政策执行。政权、军队或制度的变化须经由道路、文书、资源和地方中介才会落实，城市、乡村、边地和流动人群不可能在同一时间承受相同结果。

本条使用《宋史》卷485〈夏国传〉；《辽史》《金史》相关本纪；黑水城西夏文文书与考古报告，并标示“纪年框架可靠；该记录仅确证政权纪年连续，年内具体事项须由本年条另证”。这些材料能够钉住年序、命令、战役或局部地点，却不能直接推出人口、税收、身份和社会动机的总量。应区分诏令与实施、条约与边境生活、军报与伤亡统计，也不将官府的族群或行政称谓写成固定的社会本质。

从区域史角度看，该事件是资源、信息与权力重新连接的节点。不同地方会因市场、交通、宗教组织、家庭策略和环境条件而产生不同反应；材料的沉默限制我们对未被记录者所能作出的判断。

\eventanchor{XIA-1054-REGNAL-YEAR}{1054年·西夏以李谅祚在位第6年作为诏令、赋役与外交文书的纪年框架}账本条目记录西夏在1054年的“西夏以李谅祚在位第6年作为诏令、赋役与外交文书的纪年框架”。其背景是新岁颁历、年号和纪年是中枢与地方行政沟通的共同前提，后续影响为为同年政令、战事和地方材料提供日期锚点，不能单独推知具体政策执行。政权、军队或制度的变化须经由道路、文书、资源和地方中介才会落实，城市、乡村、边地和流动人群不可能在同一时间承受相同结果。

本条使用《宋史》卷485〈夏国传〉；《辽史》《金史》相关本纪；黑水城西夏文文书与考古报告，并标示“纪年框架可靠；该记录仅确证政权纪年连续，年内具体事项须由本年条另证”。这些材料能够钉住年序、命令、战役或局部地点，却不能直接推出人口、税收、身份和社会动机的总量。应区分诏令与实施、条约与边境生活、军报与伤亡统计，也不将官府的族群或行政称谓写成固定的社会本质。

从区域史角度看，该事件是资源、信息与权力重新连接的节点。不同地方会因市场、交通、宗教组织、家庭策略和环境条件而产生不同反应；材料的沉默限制我们对未被记录者所能作出的判断。

\eventanchor{XIA-1055-REGNAL-YEAR}{1055年·西夏以李谅祚在位第7年作为诏令、赋役与外交文书的纪年框架}账本条目记录西夏在1055年的“西夏以李谅祚在位第7年作为诏令、赋役与外交文书的纪年框架”。其背景是新岁颁历、年号和纪年是中枢与地方行政沟通的共同前提，后续影响为为同年政令、战事和地方材料提供日期锚点，不能单独推知具体政策执行。政权、军队或制度的变化须经由道路、文书、资源和地方中介才会落实，城市、乡村、边地和流动人群不可能在同一时间承受相同结果。

本条使用《宋史》卷485〈夏国传〉；《辽史》《金史》相关本纪；黑水城西夏文文书与考古报告，并标示“纪年框架可靠；该记录仅确证政权纪年连续，年内具体事项须由本年条另证”。这些材料能够钉住年序、命令、战役或局部地点，却不能直接推出人口、税收、身份和社会动机的总量。应区分诏令与实施、条约与边境生活、军报与伤亡统计，也不将官府的族群或行政称谓写成固定的社会本质。

从区域史角度看，该事件是资源、信息与权力重新连接的节点。不同地方会因市场、交通、宗教组织、家庭策略和环境条件而产生不同反应；材料的沉默限制我们对未被记录者所能作出的判断。

\eventanchor{XIA-1056-REGNAL-YEAR}{1056年·西夏以李谅祚在位第8年作为诏令、赋役与外交文书的纪年框架}账本条目记录西夏在1056年的“西夏以李谅祚在位第8年作为诏令、赋役与外交文书的纪年框架”。其背景是新岁颁历、年号和纪年是中枢与地方行政沟通的共同前提，后续影响为为同年政令、战事和地方材料提供日期锚点，不能单独推知具体政策执行。政权、军队或制度的变化须经由道路、文书、资源和地方中介才会落实，城市、乡村、边地和流动人群不可能在同一时间承受相同结果。

本条使用《宋史》卷485〈夏国传〉；《辽史》《金史》相关本纪；黑水城西夏文文书与考古报告，并标示“纪年框架可靠；该记录仅确证政权纪年连续，年内具体事项须由本年条另证”。这些材料能够钉住年序、命令、战役或局部地点，却不能直接推出人口、税收、身份和社会动机的总量。应区分诏令与实施、条约与边境生活、军报与伤亡统计，也不将官府的族群或行政称谓写成固定的社会本质。

从区域史角度看，该事件是资源、信息与权力重新连接的节点。不同地方会因市场、交通、宗教组织、家庭策略和环境条件而产生不同反应；材料的沉默限制我们对未被记录者所能作出的判断。

\eventanchor{XIA-1057-REGNAL-YEAR}{1057年·西夏以李谅祚在位第9年作为诏令、赋役与外交文书的纪年框架}账本条目记录西夏在1057年的“西夏以李谅祚在位第9年作为诏令、赋役与外交文书的纪年框架”。其背景是新岁颁历、年号和纪年是中枢与地方行政沟通的共同前提，后续影响为为同年政令、战事和地方材料提供日期锚点，不能单独推知具体政策执行。政权、军队或制度的变化须经由道路、文书、资源和地方中介才会落实，城市、乡村、边地和流动人群不可能在同一时间承受相同结果。

本条使用《宋史》卷485〈夏国传〉；《辽史》《金史》相关本纪；黑水城西夏文文书与考古报告，并标示“纪年框架可靠；该记录仅确证政权纪年连续，年内具体事项须由本年条另证”。这些材料能够钉住年序、命令、战役或局部地点，却不能直接推出人口、税收、身份和社会动机的总量。应区分诏令与实施、条约与边境生活、军报与伤亡统计，也不将官府的族群或行政称谓写成固定的社会本质。

从区域史角度看，该事件是资源、信息与权力重新连接的节点。不同地方会因市场、交通、宗教组织、家庭策略和环境条件而产生不同反应；材料的沉默限制我们对未被记录者所能作出的判断。

\eventanchor{XIA-1058-REGNAL-YEAR}{1058年·西夏以李谅祚在位第10年作为诏令、赋役与外交文书的纪年框架}账本条目记录西夏在1058年的“西夏以李谅祚在位第10年作为诏令、赋役与外交文书的纪年框架”。其背景是新岁颁历、年号和纪年是中枢与地方行政沟通的共同前提，后续影响为为同年政令、战事和地方材料提供日期锚点，不能单独推知具体政策执行。政权、军队或制度的变化须经由道路、文书、资源和地方中介才会落实，城市、乡村、边地和流动人群不可能在同一时间承受相同结果。

本条使用《宋史》卷485〈夏国传〉；《辽史》《金史》相关本纪；黑水城西夏文文书与考古报告，并标示“纪年框架可靠；该记录仅确证政权纪年连续，年内具体事项须由本年条另证”。这些材料能够钉住年序、命令、战役或局部地点，却不能直接推出人口、税收、身份和社会动机的总量。应区分诏令与实施、条约与边境生活、军报与伤亡统计，也不将官府的族群或行政称谓写成固定的社会本质。

从区域史角度看，该事件是资源、信息与权力重新连接的节点。不同地方会因市场、交通、宗教组织、家庭策略和环境条件而产生不同反应；材料的沉默限制我们对未被记录者所能作出的判断。

\eventanchor{XIA-1059-REGNAL-YEAR}{1059年·西夏以李谅祚在位第11年作为诏令、赋役与外交文书的纪年框架}账本条目记录西夏在1059年的“西夏以李谅祚在位第11年作为诏令、赋役与外交文书的纪年框架”。其背景是新岁颁历、年号和纪年是中枢与地方行政沟通的共同前提，后续影响为为同年政令、战事和地方材料提供日期锚点，不能单独推知具体政策执行。政权、军队或制度的变化须经由道路、文书、资源和地方中介才会落实，城市、乡村、边地和流动人群不可能在同一时间承受相同结果。

本条使用《宋史》卷485〈夏国传〉；《辽史》《金史》相关本纪；黑水城西夏文文书与考古报告，并标示“纪年框架可靠；该记录仅确证政权纪年连续，年内具体事项须由本年条另证”。这些材料能够钉住年序、命令、战役或局部地点，却不能直接推出人口、税收、身份和社会动机的总量。应区分诏令与实施、条约与边境生活、军报与伤亡统计，也不将官府的族群或行政称谓写成固定的社会本质。

从区域史角度看，该事件是资源、信息与权力重新连接的节点。不同地方会因市场、交通、宗教组织、家庭策略和环境条件而产生不同反应；材料的沉默限制我们对未被记录者所能作出的判断。

\eventanchor{XIA-1060-REGNAL-YEAR}{1060年·西夏以李谅祚在位第12年作为诏令、赋役与外交文书的纪年框架}账本条目记录西夏在1060年的“西夏以李谅祚在位第12年作为诏令、赋役与外交文书的纪年框架”。其背景是新岁颁历、年号和纪年是中枢与地方行政沟通的共同前提，后续影响为为同年政令、战事和地方材料提供日期锚点，不能单独推知具体政策执行。政权、军队或制度的变化须经由道路、文书、资源和地方中介才会落实，城市、乡村、边地和流动人群不可能在同一时间承受相同结果。

本条使用《宋史》卷485〈夏国传〉；《辽史》《金史》相关本纪；黑水城西夏文文书与考古报告，并标示“纪年框架可靠；该记录仅确证政权纪年连续，年内具体事项须由本年条另证”。这些材料能够钉住年序、命令、战役或局部地点，却不能直接推出人口、税收、身份和社会动机的总量。应区分诏令与实施、条约与边境生活、军报与伤亡统计，也不将官府的族群或行政称谓写成固定的社会本质。

从区域史角度看，该事件是资源、信息与权力重新连接的节点。不同地方会因市场、交通、宗教组织、家庭策略和环境条件而产生不同反应；材料的沉默限制我们对未被记录者所能作出的判断。

\eventanchor{XIA-1061-REGNAL-YEAR}{1061年·西夏以李谅祚在位第13年作为诏令、赋役与外交文书的纪年框架}账本条目记录西夏在1061年的“西夏以李谅祚在位第13年作为诏令、赋役与外交文书的纪年框架”。其背景是新岁颁历、年号和纪年是中枢与地方行政沟通的共同前提，后续影响为为同年政令、战事和地方材料提供日期锚点，不能单独推知具体政策执行。政权、军队或制度的变化须经由道路、文书、资源和地方中介才会落实，城市、乡村、边地和流动人群不可能在同一时间承受相同结果。

本条使用《宋史》卷485〈夏国传〉；《辽史》《金史》相关本纪；黑水城西夏文文书与考古报告，并标示“纪年框架可靠；该记录仅确证政权纪年连续，年内具体事项须由本年条另证”。这些材料能够钉住年序、命令、战役或局部地点，却不能直接推出人口、税收、身份和社会动机的总量。应区分诏令与实施、条约与边境生活、军报与伤亡统计，也不将官府的族群或行政称谓写成固定的社会本质。

从区域史角度看，该事件是资源、信息与权力重新连接的节点。不同地方会因市场、交通、宗教组织、家庭策略和环境条件而产生不同反应；材料的沉默限制我们对未被记录者所能作出的判断。

\eventanchor{XIA-1062-REGNAL-YEAR}{1062年·西夏以李谅祚在位第14年作为诏令、赋役与外交文书的纪年框架}账本条目记录西夏在1062年的“西夏以李谅祚在位第14年作为诏令、赋役与外交文书的纪年框架”。其背景是新岁颁历、年号和纪年是中枢与地方行政沟通的共同前提，后续影响为为同年政令、战事和地方材料提供日期锚点，不能单独推知具体政策执行。政权、军队或制度的变化须经由道路、文书、资源和地方中介才会落实，城市、乡村、边地和流动人群不可能在同一时间承受相同结果。

本条使用《宋史》卷485〈夏国传〉；《辽史》《金史》相关本纪；黑水城西夏文文书与考古报告，并标示“纪年框架可靠；该记录仅确证政权纪年连续，年内具体事项须由本年条另证”。这些材料能够钉住年序、命令、战役或局部地点，却不能直接推出人口、税收、身份和社会动机的总量。应区分诏令与实施、条约与边境生活、军报与伤亡统计，也不将官府的族群或行政称谓写成固定的社会本质。

从区域史角度看，该事件是资源、信息与权力重新连接的节点。不同地方会因市场、交通、宗教组织、家庭策略和环境条件而产生不同反应；材料的沉默限制我们对未被记录者所能作出的判断。

\eventanchor{XIA-1063-REGNAL-YEAR}{1063年·西夏以李谅祚在位第15年作为诏令、赋役与外交文书的纪年框架}账本条目记录西夏在1063年的“西夏以李谅祚在位第15年作为诏令、赋役与外交文书的纪年框架”。其背景是新岁颁历、年号和纪年是中枢与地方行政沟通的共同前提，后续影响为为同年政令、战事和地方材料提供日期锚点，不能单独推知具体政策执行。政权、军队或制度的变化须经由道路、文书、资源和地方中介才会落实，城市、乡村、边地和流动人群不可能在同一时间承受相同结果。

本条使用《宋史》卷485〈夏国传〉；《辽史》《金史》相关本纪；黑水城西夏文文书与考古报告，并标示“纪年框架可靠；该记录仅确证政权纪年连续，年内具体事项须由本年条另证”。这些材料能够钉住年序、命令、战役或局部地点，却不能直接推出人口、税收、身份和社会动机的总量。应区分诏令与实施、条约与边境生活、军报与伤亡统计，也不将官府的族群或行政称谓写成固定的社会本质。

从区域史角度看，该事件是资源、信息与权力重新连接的节点。不同地方会因市场、交通、宗教组织、家庭策略和环境条件而产生不同反应；材料的沉默限制我们对未被记录者所能作出的判断。

\eventanchor{XIA-1064-REGNAL-YEAR}{1064年·西夏以李谅祚在位第16年作为诏令、赋役与外交文书的纪年框架}账本条目记录西夏在1064年的“西夏以李谅祚在位第16年作为诏令、赋役与外交文书的纪年框架”。其背景是新岁颁历、年号和纪年是中枢与地方行政沟通的共同前提，后续影响为为同年政令、战事和地方材料提供日期锚点，不能单独推知具体政策执行。政权、军队或制度的变化须经由道路、文书、资源和地方中介才会落实，城市、乡村、边地和流动人群不可能在同一时间承受相同结果。

本条使用《宋史》卷485〈夏国传〉；《辽史》《金史》相关本纪；黑水城西夏文文书与考古报告，并标示“纪年框架可靠；该记录仅确证政权纪年连续，年内具体事项须由本年条另证”。这些材料能够钉住年序、命令、战役或局部地点，却不能直接推出人口、税收、身份和社会动机的总量。应区分诏令与实施、条约与边境生活、军报与伤亡统计，也不将官府的族群或行政称谓写成固定的社会本质。

从区域史角度看，该事件是资源、信息与权力重新连接的节点。不同地方会因市场、交通、宗教组织、家庭策略和环境条件而产生不同反应；材料的沉默限制我们对未被记录者所能作出的判断。

\eventanchor{XIA-1065-REGNAL-YEAR}{1065年·西夏以李谅祚在位第17年作为诏令、赋役与外交文书的纪年框架}账本条目记录西夏在1065年的“西夏以李谅祚在位第17年作为诏令、赋役与外交文书的纪年框架”。其背景是新岁颁历、年号和纪年是中枢与地方行政沟通的共同前提，后续影响为为同年政令、战事和地方材料提供日期锚点，不能单独推知具体政策执行。政权、军队或制度的变化须经由道路、文书、资源和地方中介才会落实，城市、乡村、边地和流动人群不可能在同一时间承受相同结果。

本条使用《宋史》卷485〈夏国传〉；《辽史》《金史》相关本纪；黑水城西夏文文书与考古报告，并标示“纪年框架可靠；该记录仅确证政权纪年连续，年内具体事项须由本年条另证”。这些材料能够钉住年序、命令、战役或局部地点，却不能直接推出人口、税收、身份和社会动机的总量。应区分诏令与实施、条约与边境生活、军报与伤亡统计，也不将官府的族群或行政称谓写成固定的社会本质。

从区域史角度看，该事件是资源、信息与权力重新连接的节点。不同地方会因市场、交通、宗教组织、家庭策略和环境条件而产生不同反应；材料的沉默限制我们对未被记录者所能作出的判断。

\eventanchor{XIA-1066-REGNAL-YEAR}{1066年·西夏以李谅祚在位第18年作为诏令、赋役与外交文书的纪年框架}账本条目记录西夏在1066年的“西夏以李谅祚在位第18年作为诏令、赋役与外交文书的纪年框架”。其背景是新岁颁历、年号和纪年是中枢与地方行政沟通的共同前提，后续影响为为同年政令、战事和地方材料提供日期锚点，不能单独推知具体政策执行。政权、军队或制度的变化须经由道路、文书、资源和地方中介才会落实，城市、乡村、边地和流动人群不可能在同一时间承受相同结果。

本条使用《宋史》卷485〈夏国传〉；《辽史》《金史》相关本纪；黑水城西夏文文书与考古报告，并标示“纪年框架可靠；该记录仅确证政权纪年连续，年内具体事项须由本年条另证”。这些材料能够钉住年序、命令、战役或局部地点，却不能直接推出人口、税收、身份和社会动机的总量。应区分诏令与实施、条约与边境生活、军报与伤亡统计，也不将官府的族群或行政称谓写成固定的社会本质。

从区域史角度看，该事件是资源、信息与权力重新连接的节点。不同地方会因市场、交通、宗教组织、家庭策略和环境条件而产生不同反应；材料的沉默限制我们对未被记录者所能作出的判断。

\eventanchor{XIA-1067-EVENT-063}{1067年·李秉常即位，西夏再度出现幼主政治}账本条目记录西夏在1067年的“李秉常即位，西夏再度出现幼主政治”。其背景是前期边防、财政和统治联盟的压力构成直接背景，具体条件须按本条材料分辨，后续影响为改变后续资源配置与政治选择；实际执行范围和社会影响仍须分项核验。政权、军队或制度的变化须经由道路、文书、资源和地方中介才会落实，城市、乡村、边地和流动人群不可能在同一时间承受相同结果。

本条使用《宋史》卷485〈夏国传〉；《辽史》《金史》相关本纪；黑水城西夏文文书与考古报告，并标示“编年主干可靠；细节、规模与动机须与对方材料、地方文书或考古材料互校”。这些材料能够钉住年序、命令、战役或局部地点，却不能直接推出人口、税收、身份和社会动机的总量。应区分诏令与实施、条约与边境生活、军报与伤亡统计，也不将官府的族群或行政称谓写成固定的社会本质。

从区域史角度看，该事件是资源、信息与权力重新连接的节点。不同地方会因市场、交通、宗教组织、家庭策略和环境条件而产生不同反应；材料的沉默限制我们对未被记录者所能作出的判断。

\eventanchor{XIA-1067-REGNAL-YEAR}{1067年·西夏以李谅祚在位第19年作为诏令、赋役与外交文书的纪年框架}账本条目记录西夏在1067年的“西夏以李谅祚在位第19年作为诏令、赋役与外交文书的纪年框架”。其背景是新岁颁历、年号和纪年是中枢与地方行政沟通的共同前提，后续影响为为同年政令、战事和地方材料提供日期锚点，不能单独推知具体政策执行。政权、军队或制度的变化须经由道路、文书、资源和地方中介才会落实，城市、乡村、边地和流动人群不可能在同一时间承受相同结果。

本条使用《宋史》卷485〈夏国传〉；《辽史》《金史》相关本纪；黑水城西夏文文书与考古报告，并标示“纪年框架可靠；该记录仅确证政权纪年连续，年内具体事项须由本年条另证”。这些材料能够钉住年序、命令、战役或局部地点，却不能直接推出人口、税收、身份和社会动机的总量。应区分诏令与实施、条约与边境生活、军报与伤亡统计，也不将官府的族群或行政称谓写成固定的社会本质。

从区域史角度看，该事件是资源、信息与权力重新连接的节点。不同地方会因市场、交通、宗教组织、家庭策略和环境条件而产生不同反应；材料的沉默限制我们对未被记录者所能作出的判断。

\eventanchor{XIA-1068-REGNAL-YEAR}{1068年·西夏以李秉常在位第1年作为诏令、赋役与外交文书的纪年框架}账本条目记录西夏在1068年的“西夏以李秉常在位第1年作为诏令、赋役与外交文书的纪年框架”。其背景是新岁颁历、年号和纪年是中枢与地方行政沟通的共同前提，后续影响为为同年政令、战事和地方材料提供日期锚点，不能单独推知具体政策执行。政权、军队或制度的变化须经由道路、文书、资源和地方中介才会落实，城市、乡村、边地和流动人群不可能在同一时间承受相同结果。

本条使用《宋史》卷485〈夏国传〉；《辽史》《金史》相关本纪；黑水城西夏文文书与考古报告，并标示“纪年框架可靠；该记录仅确证政权纪年连续，年内具体事项须由本年条另证”。这些材料能够钉住年序、命令、战役或局部地点，却不能直接推出人口、税收、身份和社会动机的总量。应区分诏令与实施、条约与边境生活、军报与伤亡统计，也不将官府的族群或行政称谓写成固定的社会本质。

从区域史角度看，该事件是资源、信息与权力重新连接的节点。不同地方会因市场、交通、宗教组织、家庭策略和环境条件而产生不同反应；材料的沉默限制我们对未被记录者所能作出的判断。

\eventanchor{XIA-1069-REGNAL-YEAR}{1069年·西夏以李秉常在位第2年作为诏令、赋役与外交文书的纪年框架}账本条目记录西夏在1069年的“西夏以李秉常在位第2年作为诏令、赋役与外交文书的纪年框架”。其背景是新岁颁历、年号和纪年是中枢与地方行政沟通的共同前提，后续影响为为同年政令、战事和地方材料提供日期锚点，不能单独推知具体政策执行。政权、军队或制度的变化须经由道路、文书、资源和地方中介才会落实，城市、乡村、边地和流动人群不可能在同一时间承受相同结果。

本条使用《宋史》卷485〈夏国传〉；《辽史》《金史》相关本纪；黑水城西夏文文书与考古报告，并标示“纪年框架可靠；该记录仅确证政权纪年连续，年内具体事项须由本年条另证”。这些材料能够钉住年序、命令、战役或局部地点，却不能直接推出人口、税收、身份和社会动机的总量。应区分诏令与实施、条约与边境生活、军报与伤亡统计，也不将官府的族群或行政称谓写成固定的社会本质。

从区域史角度看，该事件是资源、信息与权力重新连接的节点。不同地方会因市场、交通、宗教组织、家庭策略和环境条件而产生不同反应；材料的沉默限制我们对未被记录者所能作出的判断。

\eventanchor{XIA-1070-REGNAL-YEAR}{1070年·西夏以李秉常在位第3年作为诏令、赋役与外交文书的纪年框架}账本条目记录西夏在1070年的“西夏以李秉常在位第3年作为诏令、赋役与外交文书的纪年框架”。其背景是新岁颁历、年号和纪年是中枢与地方行政沟通的共同前提，后续影响为为同年政令、战事和地方材料提供日期锚点，不能单独推知具体政策执行。政权、军队或制度的变化须经由道路、文书、资源和地方中介才会落实，城市、乡村、边地和流动人群不可能在同一时间承受相同结果。

本条使用《宋史》卷485〈夏国传〉；《辽史》《金史》相关本纪；黑水城西夏文文书与考古报告，并标示“纪年框架可靠；该记录仅确证政权纪年连续，年内具体事项须由本年条另证”。这些材料能够钉住年序、命令、战役或局部地点，却不能直接推出人口、税收、身份和社会动机的总量。应区分诏令与实施、条约与边境生活、军报与伤亡统计，也不将官府的族群或行政称谓写成固定的社会本质。

从区域史角度看，该事件是资源、信息与权力重新连接的节点。不同地方会因市场、交通、宗教组织、家庭策略和环境条件而产生不同反应；材料的沉默限制我们对未被记录者所能作出的判断。

\eventanchor{XIA-1071-REGNAL-YEAR}{1071年·西夏以李秉常在位第4年作为诏令、赋役与外交文书的纪年框架}账本条目记录西夏在1071年的“西夏以李秉常在位第4年作为诏令、赋役与外交文书的纪年框架”。其背景是新岁颁历、年号和纪年是中枢与地方行政沟通的共同前提，后续影响为为同年政令、战事和地方材料提供日期锚点，不能单独推知具体政策执行。政权、军队或制度的变化须经由道路、文书、资源和地方中介才会落实，城市、乡村、边地和流动人群不可能在同一时间承受相同结果。

本条使用《宋史》卷485〈夏国传〉；《辽史》《金史》相关本纪；黑水城西夏文文书与考古报告，并标示“纪年框架可靠；该记录仅确证政权纪年连续，年内具体事项须由本年条另证”。这些材料能够钉住年序、命令、战役或局部地点，却不能直接推出人口、税收、身份和社会动机的总量。应区分诏令与实施、条约与边境生活、军报与伤亡统计，也不将官府的族群或行政称谓写成固定的社会本质。

从区域史角度看，该事件是资源、信息与权力重新连接的节点。不同地方会因市场、交通、宗教组织、家庭策略和环境条件而产生不同反应；材料的沉默限制我们对未被记录者所能作出的判断。

\eventanchor{XIA-1072-REGNAL-YEAR}{1072年·西夏以李秉常在位第5年作为诏令、赋役与外交文书的纪年框架}账本条目记录西夏在1072年的“西夏以李秉常在位第5年作为诏令、赋役与外交文书的纪年框架”。其背景是新岁颁历、年号和纪年是中枢与地方行政沟通的共同前提，后续影响为为同年政令、战事和地方材料提供日期锚点，不能单独推知具体政策执行。政权、军队或制度的变化须经由道路、文书、资源和地方中介才会落实，城市、乡村、边地和流动人群不可能在同一时间承受相同结果。

本条使用《宋史》卷485〈夏国传〉；《辽史》《金史》相关本纪；黑水城西夏文文书与考古报告，并标示“纪年框架可靠；该记录仅确证政权纪年连续，年内具体事项须由本年条另证”。这些材料能够钉住年序、命令、战役或局部地点，却不能直接推出人口、税收、身份和社会动机的总量。应区分诏令与实施、条约与边境生活、军报与伤亡统计，也不将官府的族群或行政称谓写成固定的社会本质。

从区域史角度看，该事件是资源、信息与权力重新连接的节点。不同地方会因市场、交通、宗教组织、家庭策略和环境条件而产生不同反应；材料的沉默限制我们对未被记录者所能作出的判断。

\eventanchor{XIA-1073-REGNAL-YEAR}{1073年·西夏以李秉常在位第6年作为诏令、赋役与外交文书的纪年框架}账本条目记录西夏在1073年的“西夏以李秉常在位第6年作为诏令、赋役与外交文书的纪年框架”。其背景是新岁颁历、年号和纪年是中枢与地方行政沟通的共同前提，后续影响为为同年政令、战事和地方材料提供日期锚点，不能单独推知具体政策执行。政权、军队或制度的变化须经由道路、文书、资源和地方中介才会落实，城市、乡村、边地和流动人群不可能在同一时间承受相同结果。

本条使用《宋史》卷485〈夏国传〉；《辽史》《金史》相关本纪；黑水城西夏文文书与考古报告，并标示“纪年框架可靠；该记录仅确证政权纪年连续，年内具体事项须由本年条另证”。这些材料能够钉住年序、命令、战役或局部地点，却不能直接推出人口、税收、身份和社会动机的总量。应区分诏令与实施、条约与边境生活、军报与伤亡统计，也不将官府的族群或行政称谓写成固定的社会本质。

从区域史角度看，该事件是资源、信息与权力重新连接的节点。不同地方会因市场、交通、宗教组织、家庭策略和环境条件而产生不同反应；材料的沉默限制我们对未被记录者所能作出的判断。

\eventanchor{XIA-1074-REGNAL-YEAR}{1074年·西夏以李秉常在位第7年作为诏令、赋役与外交文书的纪年框架}账本条目记录西夏在1074年的“西夏以李秉常在位第7年作为诏令、赋役与外交文书的纪年框架”。其背景是新岁颁历、年号和纪年是中枢与地方行政沟通的共同前提，后续影响为为同年政令、战事和地方材料提供日期锚点，不能单独推知具体政策执行。政权、军队或制度的变化须经由道路、文书、资源和地方中介才会落实，城市、乡村、边地和流动人群不可能在同一时间承受相同结果。

本条使用《宋史》卷485〈夏国传〉；《辽史》《金史》相关本纪；黑水城西夏文文书与考古报告，并标示“纪年框架可靠；该记录仅确证政权纪年连续，年内具体事项须由本年条另证”。这些材料能够钉住年序、命令、战役或局部地点，却不能直接推出人口、税收、身份和社会动机的总量。应区分诏令与实施、条约与边境生活、军报与伤亡统计，也不将官府的族群或行政称谓写成固定的社会本质。

从区域史角度看，该事件是资源、信息与权力重新连接的节点。不同地方会因市场、交通、宗教组织、家庭策略和环境条件而产生不同反应；材料的沉默限制我们对未被记录者所能作出的判断。

\eventanchor{XIA-1075-REGNAL-YEAR}{1075年·西夏以李秉常在位第8年作为诏令、赋役与外交文书的纪年框架}账本条目记录西夏在1075年的“西夏以李秉常在位第8年作为诏令、赋役与外交文书的纪年框架”。其背景是新岁颁历、年号和纪年是中枢与地方行政沟通的共同前提，后续影响为为同年政令、战事和地方材料提供日期锚点，不能单独推知具体政策执行。政权、军队或制度的变化须经由道路、文书、资源和地方中介才会落实，城市、乡村、边地和流动人群不可能在同一时间承受相同结果。

本条使用《宋史》卷485〈夏国传〉；《辽史》《金史》相关本纪；黑水城西夏文文书与考古报告，并标示“纪年框架可靠；该记录仅确证政权纪年连续，年内具体事项须由本年条另证”。这些材料能够钉住年序、命令、战役或局部地点，却不能直接推出人口、税收、身份和社会动机的总量。应区分诏令与实施、条约与边境生活、军报与伤亡统计，也不将官府的族群或行政称谓写成固定的社会本质。

从区域史角度看，该事件是资源、信息与权力重新连接的节点。不同地方会因市场、交通、宗教组织、家庭策略和环境条件而产生不同反应；材料的沉默限制我们对未被记录者所能作出的判断。

\eventanchor{XIA-1076-REGNAL-YEAR}{1076年·西夏以李秉常在位第9年作为诏令、赋役与外交文书的纪年框架}账本条目记录西夏在1076年的“西夏以李秉常在位第9年作为诏令、赋役与外交文书的纪年框架”。其背景是新岁颁历、年号和纪年是中枢与地方行政沟通的共同前提，后续影响为为同年政令、战事和地方材料提供日期锚点，不能单独推知具体政策执行。政权、军队或制度的变化须经由道路、文书、资源和地方中介才会落实，城市、乡村、边地和流动人群不可能在同一时间承受相同结果。

本条使用《宋史》卷485〈夏国传〉；《辽史》《金史》相关本纪；黑水城西夏文文书与考古报告，并标示“纪年框架可靠；该记录仅确证政权纪年连续，年内具体事项须由本年条另证”。这些材料能够钉住年序、命令、战役或局部地点，却不能直接推出人口、税收、身份和社会动机的总量。应区分诏令与实施、条约与边境生活、军报与伤亡统计，也不将官府的族群或行政称谓写成固定的社会本质。

从区域史角度看，该事件是资源、信息与权力重新连接的节点。不同地方会因市场、交通、宗教组织、家庭策略和环境条件而产生不同反应；材料的沉默限制我们对未被记录者所能作出的判断。

\subsection{制度、区域与材料边界}

\eventanchor{XIA-1077-REGNAL-YEAR}{1077年·西夏以李秉常在位第10年作为诏令、赋役与外交文书的纪年框架}账本条目记录西夏在1077年的“西夏以李秉常在位第10年作为诏令、赋役与外交文书的纪年框架”。其背景是新岁颁历、年号和纪年是中枢与地方行政沟通的共同前提，后续影响为为同年政令、战事和地方材料提供日期锚点，不能单独推知具体政策执行。政权、军队或制度的变化须经由道路、文书、资源和地方中介才会落实，城市、乡村、边地和流动人群不可能在同一时间承受相同结果。

本条使用《宋史》卷485〈夏国传〉；《辽史》《金史》相关本纪；黑水城西夏文文书与考古报告，并标示“纪年框架可靠；该记录仅确证政权纪年连续，年内具体事项须由本年条另证”。这些材料能够钉住年序、命令、战役或局部地点，却不能直接推出人口、税收、身份和社会动机的总量。应区分诏令与实施、条约与边境生活、军报与伤亡统计，也不将官府的族群或行政称谓写成固定的社会本质。

从区域史角度看，该事件是资源、信息与权力重新连接的节点。不同地方会因市场、交通、宗教组织、家庭策略和环境条件而产生不同反应；材料的沉默限制我们对未被记录者所能作出的判断。

\eventanchor{XIA-1078-REGNAL-YEAR}{1078年·西夏以李秉常在位第11年作为诏令、赋役与外交文书的纪年框架}账本条目记录西夏在1078年的“西夏以李秉常在位第11年作为诏令、赋役与外交文书的纪年框架”。其背景是新岁颁历、年号和纪年是中枢与地方行政沟通的共同前提，后续影响为为同年政令、战事和地方材料提供日期锚点，不能单独推知具体政策执行。政权、军队或制度的变化须经由道路、文书、资源和地方中介才会落实，城市、乡村、边地和流动人群不可能在同一时间承受相同结果。

本条使用《宋史》卷485〈夏国传〉；《辽史》《金史》相关本纪；黑水城西夏文文书与考古报告，并标示“纪年框架可靠；该记录仅确证政权纪年连续，年内具体事项须由本年条另证”。这些材料能够钉住年序、命令、战役或局部地点，却不能直接推出人口、税收、身份和社会动机的总量。应区分诏令与实施、条约与边境生活、军报与伤亡统计，也不将官府的族群或行政称谓写成固定的社会本质。

从区域史角度看，该事件是资源、信息与权力重新连接的节点。不同地方会因市场、交通、宗教组织、家庭策略和环境条件而产生不同反应；材料的沉默限制我们对未被记录者所能作出的判断。

\eventanchor{XIA-1079-REGNAL-YEAR}{1079年·西夏以李秉常在位第12年作为诏令、赋役与外交文书的纪年框架}账本条目记录西夏在1079年的“西夏以李秉常在位第12年作为诏令、赋役与外交文书的纪年框架”。其背景是新岁颁历、年号和纪年是中枢与地方行政沟通的共同前提，后续影响为为同年政令、战事和地方材料提供日期锚点，不能单独推知具体政策执行。政权、军队或制度的变化须经由道路、文书、资源和地方中介才会落实，城市、乡村、边地和流动人群不可能在同一时间承受相同结果。

本条使用《宋史》卷485〈夏国传〉；《辽史》《金史》相关本纪；黑水城西夏文文书与考古报告，并标示“纪年框架可靠；该记录仅确证政权纪年连续，年内具体事项须由本年条另证”。这些材料能够钉住年序、命令、战役或局部地点，却不能直接推出人口、税收、身份和社会动机的总量。应区分诏令与实施、条约与边境生活、军报与伤亡统计，也不将官府的族群或行政称谓写成固定的社会本质。

从区域史角度看，该事件是资源、信息与权力重新连接的节点。不同地方会因市场、交通、宗教组织、家庭策略和环境条件而产生不同反应；材料的沉默限制我们对未被记录者所能作出的判断。

\eventanchor{XIA-1080-EVENT-064}{1080年·宋神宗西北进取使宋夏边境冲突升级}账本条目记录西夏与宋在1080年的“宋神宗西北进取使宋夏边境冲突升级”。其背景是前期边防、财政和统治联盟的压力构成直接背景，具体条件须按本条材料分辨，后续影响为改变后续资源配置与政治选择；实际执行范围和社会影响仍须分项核验。政权、军队或制度的变化须经由道路、文书、资源和地方中介才会落实，城市、乡村、边地和流动人群不可能在同一时间承受相同结果。

本条使用《宋史》卷485〈夏国传〉；《辽史》《金史》相关本纪；黑水城西夏文文书与考古报告，并标示“编年主干可靠；细节、规模与动机须与对方材料、地方文书或考古材料互校”。这些材料能够钉住年序、命令、战役或局部地点，却不能直接推出人口、税收、身份和社会动机的总量。应区分诏令与实施、条约与边境生活、军报与伤亡统计，也不将官府的族群或行政称谓写成固定的社会本质。

从区域史角度看，该事件是资源、信息与权力重新连接的节点。不同地方会因市场、交通、宗教组织、家庭策略和环境条件而产生不同反应；材料的沉默限制我们对未被记录者所能作出的判断。

\eventanchor{XIA-1080-REGNAL-YEAR}{1080年·西夏以李秉常在位第13年作为诏令、赋役与外交文书的纪年框架}账本条目记录西夏在1080年的“西夏以李秉常在位第13年作为诏令、赋役与外交文书的纪年框架”。其背景是新岁颁历、年号和纪年是中枢与地方行政沟通的共同前提，后续影响为为同年政令、战事和地方材料提供日期锚点，不能单独推知具体政策执行。政权、军队或制度的变化须经由道路、文书、资源和地方中介才会落实，城市、乡村、边地和流动人群不可能在同一时间承受相同结果。

本条使用《宋史》卷485〈夏国传〉；《辽史》《金史》相关本纪；黑水城西夏文文书与考古报告，并标示“纪年框架可靠；该记录仅确证政权纪年连续，年内具体事项须由本年条另证”。这些材料能够钉住年序、命令、战役或局部地点，却不能直接推出人口、税收、身份和社会动机的总量。应区分诏令与实施、条约与边境生活、军报与伤亡统计，也不将官府的族群或行政称谓写成固定的社会本质。

从区域史角度看，该事件是资源、信息与权力重新连接的节点。不同地方会因市场、交通、宗教组织、家庭策略和环境条件而产生不同反应；材料的沉默限制我们对未被记录者所能作出的判断。

\eventanchor{XIA-1081-REGNAL-YEAR}{1081年·西夏以李秉常在位第14年作为诏令、赋役与外交文书的纪年框架}账本条目记录西夏在1081年的“西夏以李秉常在位第14年作为诏令、赋役与外交文书的纪年框架”。其背景是新岁颁历、年号和纪年是中枢与地方行政沟通的共同前提，后续影响为为同年政令、战事和地方材料提供日期锚点，不能单独推知具体政策执行。政权、军队或制度的变化须经由道路、文书、资源和地方中介才会落实，城市、乡村、边地和流动人群不可能在同一时间承受相同结果。

本条使用《宋史》卷485〈夏国传〉；《辽史》《金史》相关本纪；黑水城西夏文文书与考古报告，并标示“纪年框架可靠；该记录仅确证政权纪年连续，年内具体事项须由本年条另证”。这些材料能够钉住年序、命令、战役或局部地点，却不能直接推出人口、税收、身份和社会动机的总量。应区分诏令与实施、条约与边境生活、军报与伤亡统计，也不将官府的族群或行政称谓写成固定的社会本质。

从区域史角度看，该事件是资源、信息与权力重新连接的节点。不同地方会因市场、交通、宗教组织、家庭策略和环境条件而产生不同反应；材料的沉默限制我们对未被记录者所能作出的判断。

\eventanchor{XIA-1082-EVENT-065}{1082年·西夏攻破永乐城}账本条目记录西夏与宋在1082年的“西夏攻破永乐城”。其背景是前期边防、财政和统治联盟的压力构成直接背景，具体条件须按本条材料分辨，后续影响为改变后续资源配置与政治选择；实际执行范围和社会影响仍须分项核验。政权、军队或制度的变化须经由道路、文书、资源和地方中介才会落实，城市、乡村、边地和流动人群不可能在同一时间承受相同结果。

本条使用《宋史》卷485〈夏国传〉；《辽史》《金史》相关本纪；黑水城西夏文文书与考古报告，并标示“编年主干可靠；细节、规模与动机须与对方材料、地方文书或考古材料互校”。这些材料能够钉住年序、命令、战役或局部地点，却不能直接推出人口、税收、身份和社会动机的总量。应区分诏令与实施、条约与边境生活、军报与伤亡统计，也不将官府的族群或行政称谓写成固定的社会本质。

从区域史角度看，该事件是资源、信息与权力重新连接的节点。不同地方会因市场、交通、宗教组织、家庭策略和环境条件而产生不同反应；材料的沉默限制我们对未被记录者所能作出的判断。

\eventanchor{XIA-1082-REGNAL-YEAR}{1082年·西夏以李秉常在位第15年作为诏令、赋役与外交文书的纪年框架}账本条目记录西夏在1082年的“西夏以李秉常在位第15年作为诏令、赋役与外交文书的纪年框架”。其背景是新岁颁历、年号和纪年是中枢与地方行政沟通的共同前提，后续影响为为同年政令、战事和地方材料提供日期锚点，不能单独推知具体政策执行。政权、军队或制度的变化须经由道路、文书、资源和地方中介才会落实，城市、乡村、边地和流动人群不可能在同一时间承受相同结果。

本条使用《宋史》卷485〈夏国传〉；《辽史》《金史》相关本纪；黑水城西夏文文书与考古报告，并标示“纪年框架可靠；该记录仅确证政权纪年连续，年内具体事项须由本年条另证”。这些材料能够钉住年序、命令、战役或局部地点，却不能直接推出人口、税收、身份和社会动机的总量。应区分诏令与实施、条约与边境生活、军报与伤亡统计，也不将官府的族群或行政称谓写成固定的社会本质。

从区域史角度看，该事件是资源、信息与权力重新连接的节点。不同地方会因市场、交通、宗教组织、家庭策略和环境条件而产生不同反应；材料的沉默限制我们对未被记录者所能作出的判断。

\eventanchor{XIA-1083-REGNAL-YEAR}{1083年·西夏以李秉常在位第16年作为诏令、赋役与外交文书的纪年框架}账本条目记录西夏在1083年的“西夏以李秉常在位第16年作为诏令、赋役与外交文书的纪年框架”。其背景是新岁颁历、年号和纪年是中枢与地方行政沟通的共同前提，后续影响为为同年政令、战事和地方材料提供日期锚点，不能单独推知具体政策执行。政权、军队或制度的变化须经由道路、文书、资源和地方中介才会落实，城市、乡村、边地和流动人群不可能在同一时间承受相同结果。

本条使用《宋史》卷485〈夏国传〉；《辽史》《金史》相关本纪；黑水城西夏文文书与考古报告，并标示“纪年框架可靠；该记录仅确证政权纪年连续，年内具体事项须由本年条另证”。这些材料能够钉住年序、命令、战役或局部地点，却不能直接推出人口、税收、身份和社会动机的总量。应区分诏令与实施、条约与边境生活、军报与伤亡统计，也不将官府的族群或行政称谓写成固定的社会本质。

从区域史角度看，该事件是资源、信息与权力重新连接的节点。不同地方会因市场、交通、宗教组织、家庭策略和环境条件而产生不同反应；材料的沉默限制我们对未被记录者所能作出的判断。

\eventanchor{XIA-1084-REGNAL-YEAR}{1084年·西夏以李秉常在位第17年作为诏令、赋役与外交文书的纪年框架}账本条目记录西夏在1084年的“西夏以李秉常在位第17年作为诏令、赋役与外交文书的纪年框架”。其背景是新岁颁历、年号和纪年是中枢与地方行政沟通的共同前提，后续影响为为同年政令、战事和地方材料提供日期锚点，不能单独推知具体政策执行。政权、军队或制度的变化须经由道路、文书、资源和地方中介才会落实，城市、乡村、边地和流动人群不可能在同一时间承受相同结果。

本条使用《宋史》卷485〈夏国传〉；《辽史》《金史》相关本纪；黑水城西夏文文书与考古报告，并标示“纪年框架可靠；该记录仅确证政权纪年连续，年内具体事项须由本年条另证”。这些材料能够钉住年序、命令、战役或局部地点，却不能直接推出人口、税收、身份和社会动机的总量。应区分诏令与实施、条约与边境生活、军报与伤亡统计，也不将官府的族群或行政称谓写成固定的社会本质。

从区域史角度看，该事件是资源、信息与权力重新连接的节点。不同地方会因市场、交通、宗教组织、家庭策略和环境条件而产生不同反应；材料的沉默限制我们对未被记录者所能作出的判断。

\eventanchor{XIA-1085-REGNAL-YEAR}{1085年·西夏以李秉常在位第18年作为诏令、赋役与外交文书的纪年框架}账本条目记录西夏在1085年的“西夏以李秉常在位第18年作为诏令、赋役与外交文书的纪年框架”。其背景是新岁颁历、年号和纪年是中枢与地方行政沟通的共同前提，后续影响为为同年政令、战事和地方材料提供日期锚点，不能单独推知具体政策执行。政权、军队或制度的变化须经由道路、文书、资源和地方中介才会落实，城市、乡村、边地和流动人群不可能在同一时间承受相同结果。

本条使用《宋史》卷485〈夏国传〉；《辽史》《金史》相关本纪；黑水城西夏文文书与考古报告，并标示“纪年框架可靠；该记录仅确证政权纪年连续，年内具体事项须由本年条另证”。这些材料能够钉住年序、命令、战役或局部地点，却不能直接推出人口、税收、身份和社会动机的总量。应区分诏令与实施、条约与边境生活、军报与伤亡统计，也不将官府的族群或行政称谓写成固定的社会本质。

从区域史角度看，该事件是资源、信息与权力重新连接的节点。不同地方会因市场、交通、宗教组织、家庭策略和环境条件而产生不同反应；材料的沉默限制我们对未被记录者所能作出的判断。

\eventanchor{XIA-1086-EVENT-066}{1086年·李乾顺即位}账本条目记录西夏在1086年的“李乾顺即位”。其背景是前期边防、财政和统治联盟的压力构成直接背景，具体条件须按本条材料分辨，后续影响为改变后续资源配置与政治选择；实际执行范围和社会影响仍须分项核验。政权、军队或制度的变化须经由道路、文书、资源和地方中介才会落实，城市、乡村、边地和流动人群不可能在同一时间承受相同结果。

本条使用《宋史》卷485〈夏国传〉；《辽史》《金史》相关本纪；黑水城西夏文文书与考古报告，并标示“编年主干可靠；细节、规模与动机须与对方材料、地方文书或考古材料互校”。这些材料能够钉住年序、命令、战役或局部地点，却不能直接推出人口、税收、身份和社会动机的总量。应区分诏令与实施、条约与边境生活、军报与伤亡统计，也不将官府的族群或行政称谓写成固定的社会本质。

从区域史角度看，该事件是资源、信息与权力重新连接的节点。不同地方会因市场、交通、宗教组织、家庭策略和环境条件而产生不同反应；材料的沉默限制我们对未被记录者所能作出的判断。

\eventanchor{XIA-1086-REGNAL-YEAR}{1086年·西夏以李秉常在位第19年作为诏令、赋役与外交文书的纪年框架}账本条目记录西夏在1086年的“西夏以李秉常在位第19年作为诏令、赋役与外交文书的纪年框架”。其背景是新岁颁历、年号和纪年是中枢与地方行政沟通的共同前提，后续影响为为同年政令、战事和地方材料提供日期锚点，不能单独推知具体政策执行。政权、军队或制度的变化须经由道路、文书、资源和地方中介才会落实，城市、乡村、边地和流动人群不可能在同一时间承受相同结果。

本条使用《宋史》卷485〈夏国传〉；《辽史》《金史》相关本纪；黑水城西夏文文书与考古报告，并标示“纪年框架可靠；该记录仅确证政权纪年连续，年内具体事项须由本年条另证”。这些材料能够钉住年序、命令、战役或局部地点，却不能直接推出人口、税收、身份和社会动机的总量。应区分诏令与实施、条约与边境生活、军报与伤亡统计，也不将官府的族群或行政称谓写成固定的社会本质。

从区域史角度看，该事件是资源、信息与权力重新连接的节点。不同地方会因市场、交通、宗教组织、家庭策略和环境条件而产生不同反应；材料的沉默限制我们对未被记录者所能作出的判断。

\eventanchor{XIA-1087-REGNAL-YEAR}{1087年·西夏以李乾顺在位第1年作为诏令、赋役与外交文书的纪年框架}账本条目记录西夏在1087年的“西夏以李乾顺在位第1年作为诏令、赋役与外交文书的纪年框架”。其背景是新岁颁历、年号和纪年是中枢与地方行政沟通的共同前提，后续影响为为同年政令、战事和地方材料提供日期锚点，不能单独推知具体政策执行。政权、军队或制度的变化须经由道路、文书、资源和地方中介才会落实，城市、乡村、边地和流动人群不可能在同一时间承受相同结果。

本条使用《宋史》卷485〈夏国传〉；《辽史》《金史》相关本纪；黑水城西夏文文书与考古报告，并标示“纪年框架可靠；该记录仅确证政权纪年连续，年内具体事项须由本年条另证”。这些材料能够钉住年序、命令、战役或局部地点，却不能直接推出人口、税收、身份和社会动机的总量。应区分诏令与实施、条约与边境生活、军报与伤亡统计，也不将官府的族群或行政称谓写成固定的社会本质。

从区域史角度看，该事件是资源、信息与权力重新连接的节点。不同地方会因市场、交通、宗教组织、家庭策略和环境条件而产生不同反应；材料的沉默限制我们对未被记录者所能作出的判断。

\eventanchor{XIA-1088-REGNAL-YEAR}{1088年·西夏以李乾顺在位第2年作为诏令、赋役与外交文书的纪年框架}账本条目记录西夏在1088年的“西夏以李乾顺在位第2年作为诏令、赋役与外交文书的纪年框架”。其背景是新岁颁历、年号和纪年是中枢与地方行政沟通的共同前提，后续影响为为同年政令、战事和地方材料提供日期锚点，不能单独推知具体政策执行。政权、军队或制度的变化须经由道路、文书、资源和地方中介才会落实，城市、乡村、边地和流动人群不可能在同一时间承受相同结果。

本条使用《宋史》卷485〈夏国传〉；《辽史》《金史》相关本纪；黑水城西夏文文书与考古报告，并标示“纪年框架可靠；该记录仅确证政权纪年连续，年内具体事项须由本年条另证”。这些材料能够钉住年序、命令、战役或局部地点，却不能直接推出人口、税收、身份和社会动机的总量。应区分诏令与实施、条约与边境生活、军报与伤亡统计，也不将官府的族群或行政称谓写成固定的社会本质。

从区域史角度看，该事件是资源、信息与权力重新连接的节点。不同地方会因市场、交通、宗教组织、家庭策略和环境条件而产生不同反应；材料的沉默限制我们对未被记录者所能作出的判断。

\eventanchor{XIA-1089-REGNAL-YEAR}{1089年·西夏以李乾顺在位第3年作为诏令、赋役与外交文书的纪年框架}账本条目记录西夏在1089年的“西夏以李乾顺在位第3年作为诏令、赋役与外交文书的纪年框架”。其背景是新岁颁历、年号和纪年是中枢与地方行政沟通的共同前提，后续影响为为同年政令、战事和地方材料提供日期锚点，不能单独推知具体政策执行。政权、军队或制度的变化须经由道路、文书、资源和地方中介才会落实，城市、乡村、边地和流动人群不可能在同一时间承受相同结果。

本条使用《宋史》卷485〈夏国传〉；《辽史》《金史》相关本纪；黑水城西夏文文书与考古报告，并标示“纪年框架可靠；该记录仅确证政权纪年连续，年内具体事项须由本年条另证”。这些材料能够钉住年序、命令、战役或局部地点，却不能直接推出人口、税收、身份和社会动机的总量。应区分诏令与实施、条约与边境生活、军报与伤亡统计，也不将官府的族群或行政称谓写成固定的社会本质。

从区域史角度看，该事件是资源、信息与权力重新连接的节点。不同地方会因市场、交通、宗教组织、家庭策略和环境条件而产生不同反应；材料的沉默限制我们对未被记录者所能作出的判断。

\eventanchor{XIA-1090-REGNAL-YEAR}{1090年·西夏以李乾顺在位第4年作为诏令、赋役与外交文书的纪年框架}账本条目记录西夏在1090年的“西夏以李乾顺在位第4年作为诏令、赋役与外交文书的纪年框架”。其背景是新岁颁历、年号和纪年是中枢与地方行政沟通的共同前提，后续影响为为同年政令、战事和地方材料提供日期锚点，不能单独推知具体政策执行。政权、军队或制度的变化须经由道路、文书、资源和地方中介才会落实，城市、乡村、边地和流动人群不可能在同一时间承受相同结果。

本条使用《宋史》卷485〈夏国传〉；《辽史》《金史》相关本纪；黑水城西夏文文书与考古报告，并标示“纪年框架可靠；该记录仅确证政权纪年连续，年内具体事项须由本年条另证”。这些材料能够钉住年序、命令、战役或局部地点，却不能直接推出人口、税收、身份和社会动机的总量。应区分诏令与实施、条约与边境生活、军报与伤亡统计，也不将官府的族群或行政称谓写成固定的社会本质。

从区域史角度看，该事件是资源、信息与权力重新连接的节点。不同地方会因市场、交通、宗教组织、家庭策略和环境条件而产生不同反应；材料的沉默限制我们对未被记录者所能作出的判断。

\eventanchor{XIA-1091-REGNAL-YEAR}{1091年·西夏以李乾顺在位第5年作为诏令、赋役与外交文书的纪年框架}账本条目记录西夏在1091年的“西夏以李乾顺在位第5年作为诏令、赋役与外交文书的纪年框架”。其背景是新岁颁历、年号和纪年是中枢与地方行政沟通的共同前提，后续影响为为同年政令、战事和地方材料提供日期锚点，不能单独推知具体政策执行。政权、军队或制度的变化须经由道路、文书、资源和地方中介才会落实，城市、乡村、边地和流动人群不可能在同一时间承受相同结果。

本条使用《宋史》卷485〈夏国传〉；《辽史》《金史》相关本纪；黑水城西夏文文书与考古报告，并标示“纪年框架可靠；该记录仅确证政权纪年连续，年内具体事项须由本年条另证”。这些材料能够钉住年序、命令、战役或局部地点，却不能直接推出人口、税收、身份和社会动机的总量。应区分诏令与实施、条约与边境生活、军报与伤亡统计，也不将官府的族群或行政称谓写成固定的社会本质。

从区域史角度看，该事件是资源、信息与权力重新连接的节点。不同地方会因市场、交通、宗教组织、家庭策略和环境条件而产生不同反应；材料的沉默限制我们对未被记录者所能作出的判断。

\eventanchor{XIA-1092-REGNAL-YEAR}{1092年·西夏以李乾顺在位第6年作为诏令、赋役与外交文书的纪年框架}账本条目记录西夏在1092年的“西夏以李乾顺在位第6年作为诏令、赋役与外交文书的纪年框架”。其背景是新岁颁历、年号和纪年是中枢与地方行政沟通的共同前提，后续影响为为同年政令、战事和地方材料提供日期锚点，不能单独推知具体政策执行。政权、军队或制度的变化须经由道路、文书、资源和地方中介才会落实，城市、乡村、边地和流动人群不可能在同一时间承受相同结果。

本条使用《宋史》卷485〈夏国传〉；《辽史》《金史》相关本纪；黑水城西夏文文书与考古报告，并标示“纪年框架可靠；该记录仅确证政权纪年连续，年内具体事项须由本年条另证”。这些材料能够钉住年序、命令、战役或局部地点，却不能直接推出人口、税收、身份和社会动机的总量。应区分诏令与实施、条约与边境生活、军报与伤亡统计，也不将官府的族群或行政称谓写成固定的社会本质。

从区域史角度看，该事件是资源、信息与权力重新连接的节点。不同地方会因市场、交通、宗教组织、家庭策略和环境条件而产生不同反应；材料的沉默限制我们对未被记录者所能作出的判断。

\eventanchor{XIA-1093-REGNAL-YEAR}{1093年·西夏以李乾顺在位第7年作为诏令、赋役与外交文书的纪年框架}账本条目记录西夏在1093年的“西夏以李乾顺在位第7年作为诏令、赋役与外交文书的纪年框架”。其背景是新岁颁历、年号和纪年是中枢与地方行政沟通的共同前提，后续影响为为同年政令、战事和地方材料提供日期锚点，不能单独推知具体政策执行。政权、军队或制度的变化须经由道路、文书、资源和地方中介才会落实，城市、乡村、边地和流动人群不可能在同一时间承受相同结果。

本条使用《宋史》卷485〈夏国传〉；《辽史》《金史》相关本纪；黑水城西夏文文书与考古报告，并标示“纪年框架可靠；该记录仅确证政权纪年连续，年内具体事项须由本年条另证”。这些材料能够钉住年序、命令、战役或局部地点，却不能直接推出人口、税收、身份和社会动机的总量。应区分诏令与实施、条约与边境生活、军报与伤亡统计，也不将官府的族群或行政称谓写成固定的社会本质。

从区域史角度看，该事件是资源、信息与权力重新连接的节点。不同地方会因市场、交通、宗教组织、家庭策略和环境条件而产生不同反应；材料的沉默限制我们对未被记录者所能作出的判断。

\eventanchor{XIA-1094-REGNAL-YEAR}{1094年·西夏以李乾顺在位第8年作为诏令、赋役与外交文书的纪年框架}账本条目记录西夏在1094年的“西夏以李乾顺在位第8年作为诏令、赋役与外交文书的纪年框架”。其背景是新岁颁历、年号和纪年是中枢与地方行政沟通的共同前提，后续影响为为同年政令、战事和地方材料提供日期锚点，不能单独推知具体政策执行。政权、军队或制度的变化须经由道路、文书、资源和地方中介才会落实，城市、乡村、边地和流动人群不可能在同一时间承受相同结果。

本条使用《宋史》卷485〈夏国传〉；《辽史》《金史》相关本纪；黑水城西夏文文书与考古报告，并标示“纪年框架可靠；该记录仅确证政权纪年连续，年内具体事项须由本年条另证”。这些材料能够钉住年序、命令、战役或局部地点，却不能直接推出人口、税收、身份和社会动机的总量。应区分诏令与实施、条约与边境生活、军报与伤亡统计，也不将官府的族群或行政称谓写成固定的社会本质。

从区域史角度看，该事件是资源、信息与权力重新连接的节点。不同地方会因市场、交通、宗教组织、家庭策略和环境条件而产生不同反应；材料的沉默限制我们对未被记录者所能作出的判断。

\eventanchor{XIA-1095-REGNAL-YEAR}{1095年·西夏以李乾顺在位第9年作为诏令、赋役与外交文书的纪年框架}账本条目记录西夏在1095年的“西夏以李乾顺在位第9年作为诏令、赋役与外交文书的纪年框架”。其背景是新岁颁历、年号和纪年是中枢与地方行政沟通的共同前提，后续影响为为同年政令、战事和地方材料提供日期锚点，不能单独推知具体政策执行。政权、军队或制度的变化须经由道路、文书、资源和地方中介才会落实，城市、乡村、边地和流动人群不可能在同一时间承受相同结果。

本条使用《宋史》卷485〈夏国传〉；《辽史》《金史》相关本纪；黑水城西夏文文书与考古报告，并标示“纪年框架可靠；该记录仅确证政权纪年连续，年内具体事项须由本年条另证”。这些材料能够钉住年序、命令、战役或局部地点，却不能直接推出人口、税收、身份和社会动机的总量。应区分诏令与实施、条约与边境生活、军报与伤亡统计，也不将官府的族群或行政称谓写成固定的社会本质。

从区域史角度看，该事件是资源、信息与权力重新连接的节点。不同地方会因市场、交通、宗教组织、家庭策略和环境条件而产生不同反应；材料的沉默限制我们对未被记录者所能作出的判断。

\eventanchor{XIA-1096-REGNAL-YEAR}{1096年·西夏以李乾顺在位第10年作为诏令、赋役与外交文书的纪年框架}账本条目记录西夏在1096年的“西夏以李乾顺在位第10年作为诏令、赋役与外交文书的纪年框架”。其背景是新岁颁历、年号和纪年是中枢与地方行政沟通的共同前提，后续影响为为同年政令、战事和地方材料提供日期锚点，不能单独推知具体政策执行。政权、军队或制度的变化须经由道路、文书、资源和地方中介才会落实，城市、乡村、边地和流动人群不可能在同一时间承受相同结果。

本条使用《宋史》卷485〈夏国传〉；《辽史》《金史》相关本纪；黑水城西夏文文书与考古报告，并标示“纪年框架可靠；该记录仅确证政权纪年连续，年内具体事项须由本年条另证”。这些材料能够钉住年序、命令、战役或局部地点，却不能直接推出人口、税收、身份和社会动机的总量。应区分诏令与实施、条约与边境生活、军报与伤亡统计，也不将官府的族群或行政称谓写成固定的社会本质。

从区域史角度看，该事件是资源、信息与权力重新连接的节点。不同地方会因市场、交通、宗教组织、家庭策略和环境条件而产生不同反应；材料的沉默限制我们对未被记录者所能作出的判断。

\eventanchor{XIA-1097-REGNAL-YEAR}{1097年·西夏以李乾顺在位第11年作为诏令、赋役与外交文书的纪年框架}账本条目记录西夏在1097年的“西夏以李乾顺在位第11年作为诏令、赋役与外交文书的纪年框架”。其背景是新岁颁历、年号和纪年是中枢与地方行政沟通的共同前提，后续影响为为同年政令、战事和地方材料提供日期锚点，不能单独推知具体政策执行。政权、军队或制度的变化须经由道路、文书、资源和地方中介才会落实，城市、乡村、边地和流动人群不可能在同一时间承受相同结果。

本条使用《宋史》卷485〈夏国传〉；《辽史》《金史》相关本纪；黑水城西夏文文书与考古报告，并标示“纪年框架可靠；该记录仅确证政权纪年连续，年内具体事项须由本年条另证”。这些材料能够钉住年序、命令、战役或局部地点，却不能直接推出人口、税收、身份和社会动机的总量。应区分诏令与实施、条约与边境生活、军报与伤亡统计，也不将官府的族群或行政称谓写成固定的社会本质。

从区域史角度看，该事件是资源、信息与权力重新连接的节点。不同地方会因市场、交通、宗教组织、家庭策略和环境条件而产生不同反应；材料的沉默限制我们对未被记录者所能作出的判断。

\eventanchor{XIA-1098-REGNAL-YEAR}{1098年·西夏以李乾顺在位第12年作为诏令、赋役与外交文书的纪年框架}账本条目记录西夏在1098年的“西夏以李乾顺在位第12年作为诏令、赋役与外交文书的纪年框架”。其背景是新岁颁历、年号和纪年是中枢与地方行政沟通的共同前提，后续影响为为同年政令、战事和地方材料提供日期锚点，不能单独推知具体政策执行。政权、军队或制度的变化须经由道路、文书、资源和地方中介才会落实，城市、乡村、边地和流动人群不可能在同一时间承受相同结果。

本条使用《宋史》卷485〈夏国传〉；《辽史》《金史》相关本纪；黑水城西夏文文书与考古报告，并标示“纪年框架可靠；该记录仅确证政权纪年连续，年内具体事项须由本年条另证”。这些材料能够钉住年序、命令、战役或局部地点，却不能直接推出人口、税收、身份和社会动机的总量。应区分诏令与实施、条约与边境生活、军报与伤亡统计，也不将官府的族群或行政称谓写成固定的社会本质。

从区域史角度看，该事件是资源、信息与权力重新连接的节点。不同地方会因市场、交通、宗教组织、家庭策略和环境条件而产生不同反应；材料的沉默限制我们对未被记录者所能作出的判断。
